\documentclass[aspectratio=169,11pt]{beamer}

\usetheme{Madrid}
\usecolortheme{whale}
\setbeamertemplate{navigation symbols}{}
\setbeamertemplate{footline}[frame number]

\usepackage[utf8]{inputenc}
\usepackage[T1]{fontenc}
\usepackage{graphicx}
\usepackage{booktabs}
\usepackage{amsmath}
\usepackage{hyperref}
\usepackage{tikz}
\usetikzlibrary{shapes,arrows,positioning,calc,fit,backgrounds}
\usepackage{xcolor}
\usepackage{array}
\usepackage{listings}

\lstset{
    basicstyle=\ttfamily\scriptsize,
    breaklines=true,
    frame=single,
    backgroundcolor=\color{gray!10},
    keywordstyle=\color{blue!70},
    commentstyle=\color{green!50!black},
    stringstyle=\color{red!60},
    showstringspaces=false,
    tabsize=2
}

\definecolor{teachbg}{RGB}{255,243,224}
\definecolor{teachborder}{RGB}{255,152,0}

\title[EECE5155 -- MLOps Exercises]{Practical Exercises: MLOps for IoT}
\subtitle{UCI HAR Dataset --- Instructor Version with Solutions}
\author{Eduardo Baena}
\date{Spring 2026}
\institute{Northeastern University\\
\texttt{e.baena@northeastern.edu}}

\begin{document}

\begin{frame}
    \titlepage
    \vspace{-1cm}
    \begin{center}
        \textbf{EECE5155: Wireless Sensor Networks and the Internet of Things}\\[4pt]
        {\small \textbf{Instructor version} --- full solutions, teaching notes, grading rubric.}
    \end{center}
\end{frame}

%==============================================================================
\begin{frame}{Exercise Overview}
    \textbf{Dataset: UCI HAR (Human Activity Recognition Using Smartphones)}\\
    {\small Anguita et al., ESANN 2013 --- \texttt{archive.ics.uci.edu/dataset/240}}

    \vspace{0.2cm}
    \begin{block}{What the dataset is}
        30 volunteers wore a Samsung Galaxy S II on the waist. 3-axis accelerometer + gyroscope at 50 Hz. Six activities: WALKING, WALKING\_UPSTAIRS, WALKING\_DOWNSTAIRS, SITTING, STANDING, LAYING. Each 2.56-second window labeled by video review. 561 statistical features per window. 10,299 total samples. Official subject-based train/test split included.
    \end{block}

    \vspace{0.2cm}
    \textbf{Why this dataset for MLOps teaching:}
    \begin{enumerate}
        \item \textbf{Real IoT hardware} --- smartphone sensors every student already owns
        \item \textbf{Explicit labeling cost} --- video review by researchers; students grasp why labels are expensive
        \item \textbf{Natural drift} --- new user populations (elderly, different phone placement) are the deployment reality
        \item \textbf{Algorithm tradeoffs are concrete} --- kNN cannot deploy on MCU, SVM can; measurable numbers
        \item \textbf{One pip install} --- no Kaggle account, no registration, no synthetic data
    \end{enumerate}
\end{frame}

%==============================================================================
\section{Dataset and Setup}
%==============================================================================

\begin{frame}[fragile]{Setup: Download Script}

\begin{lstlisting}[language=Python]
import urllib.request, zipfile, io, pandas as pd

URL = ("https://archive.ics.uci.edu/static/public/240/"
       "human+activity+recognition+using+smartphones.zip")
req = urllib.request.Request(URL, headers={"User-Agent": "Mozilla/5.0"})
with urllib.request.urlopen(req, timeout=120) as r:
    outer = zipfile.ZipFile(io.BytesIO(r.read()))
inner = zipfile.ZipFile(io.BytesIO(outer.read("UCI HAR Dataset.zip")))

def read(path):
    with inner.open(path) as f: return f.read().decode()

features = [l.split(" ",1)[1].strip()
            for l in read("UCI HAR Dataset/features.txt").strip().split("\n")]
labels   = {int(l.split()[0]): l.split()[1]
            for l in read("UCI HAR Dataset/activity_labels.txt").strip().split("\n")}

def load(split):
    X = pd.read_csv(io.StringIO(read(f"UCI HAR Dataset/{split}/X_{split}.txt")),
                    sep=r"\s+", header=None, names=features)
    y = pd.read_csv(io.StringIO(read(f"UCI HAR Dataset/{split}/y_{split}.txt")),
                    header=None)[0].map(labels).rename("Activity")
    return X, y

X_train, y_train = load("train")   # 7352 samples, subjects 1-21
X_test,  y_test  = load("test")    # 2947 samples, subjects 22-30
X_train.to_csv("har_X_train.csv", index=False)
X_test.to_csv( "har_X_test.csv",  index=False)
y_train.to_csv("har_y_train.csv", index=False)
y_test.to_csv( "har_y_test.csv",  index=False)
print(f"Train: {X_train.shape}   Test: {X_test.shape}")
\end{lstlisting}

    \begin{alertblock}{Teaching note: download approach}
        \texttt{ucimlrepo.fetch\_ucirepo(id=240)} is listed as ``not available for import'' by the UCI API (verified 2025). The direct zip download above always works and produces identical data. Run \texttt{download\_har.py} from the seminar repo once before class.
    \end{alertblock}
\end{frame}

%==============================================================================
\section{Exercise 1: Training Data and the Evaluation Trap}
%==============================================================================

\begin{frame}{Exercise 1: What Is Training Data?}
    \begin{block}{Scenario}
        A hospital wants to deploy a wearable activity monitor for post-surgery recovery patients. You have UCI HAR data from healthy young adults ages 19--48. Train a classifier and understand what the training data represents before deploying.
    \end{block}

    \vspace{0.2cm}
    \textbf{Tasks:}
    \begin{enumerate}
        \item Plot the class balance. Is the dataset balanced? Does imbalance affect metric choice?
        \item Plot a 2D scatter of two features for all 6 classes. Are some activities linearly separable?
        \item Train a Random Forest on the official train split. Print accuracy and classification report.
        \item Try random \texttt{train\_test\_split} on the same data. Compare the two accuracy numbers and explain the gap.
    \end{enumerate}

    \vspace{0.1cm}
    \textbf{Expected results:} Subject split: 0.929. Random split: $\sim$0.977. The 5\% gap is subject leakage.
\end{frame}

\begin{frame}[fragile]{Exercise 1: Solution --- Training and Split Comparison}

\begin{lstlisting}[language=Python]
import pandas as pd
from sklearn.ensemble import RandomForestClassifier
from sklearn.model_selection import train_test_split
from sklearn.metrics import accuracy_score, classification_report
import pickle

X_train = pd.read_csv("har_X_train.csv")
X_test  = pd.read_csv("har_X_test.csv")
y_train = pd.read_csv("har_y_train.csv").squeeze()
y_test  = pd.read_csv("har_y_test.csv").squeeze()

rf = RandomForestClassifier(n_estimators=100, random_state=42, n_jobs=-1)
rf.fit(X_train, y_train)
acc_correct = accuracy_score(y_test, rf.predict(X_test))
print(f"Subject split accuracy:  {acc_correct:.3f}")
print(classification_report(y_test, rf.predict(X_test)))

# Wrong split
X_all = pd.concat([X_train, X_test], ignore_index=True)
y_all = pd.concat([y_train, y_test], ignore_index=True)
X_tr2, X_te2, y_tr2, y_te2 = train_test_split(X_all, y_all, test_size=0.2, random_state=42)
rf2 = RandomForestClassifier(n_estimators=100, random_state=42, n_jobs=-1)
rf2.fit(X_tr2, y_tr2)
print(f"Random split accuracy:   {accuracy_score(y_te2, rf2.predict(X_te2)):.3f}")
pickle.dump(rf, open("rf_model.pkl", "wb"))
\end{lstlisting}

    \begin{exampleblock}{Why the 5\% gap}
        Random split leaks windows from the same subjects into both sets. The model memorizes per-person gait patterns. On new users (the actual deployment scenario), real accuracy is 0.93 not 0.98. \textbf{Evaluation protocol must match deployment scenario.} This is one of the most common ML mistakes in production IoT systems.
    \end{exampleblock}
\end{frame}

\begin{frame}[fragile]{Exercise 1: Solution --- Visualization}

\begin{lstlisting}[language=Python]
import matplotlib.pyplot as plt, seaborn as sns
from sklearn.metrics import ConfusionMatrixDisplay

# Class balance
y.value_counts().sort_values().plot(kind="barh",
    title="UCI HAR: Samples per activity", color="steelblue")
plt.tight_layout(); plt.savefig("class_balance.png"); plt.close()

# Feature scatter: features with real class separation
X_all = pd.concat([X_train, X_test], ignore_index=True)
y_all = pd.concat([y_train, y_test], ignore_index=True)
fig, ax = plt.subplots(figsize=(9, 4))
palette = sns.color_palette("tab10", 6)
for i, act in enumerate(y_all.unique()):
    mask = (y_all == act)
    ax.scatter(X_all.loc[mask, "tBodyAccMag-mean()"],
               X_all.loc[mask, "tGravityAcc-mean()-X"],
               label=act, alpha=0.25, s=8, color=palette[i])
ax.set_xlabel("Body Acc Magnitude (mean)");  ax.set_ylabel("Gravity Acc X (mean)")
ax.legend(fontsize=7, ncol=2, markerscale=2)
plt.tight_layout(); plt.savefig("feature_scatter.png"); plt.close()

# Confusion matrix (normalised by row = per-class recall)
ConfusionMatrixDisplay.from_estimator(
    rf, X_test, y_test, normalize="true",
    xticks_rotation=45, colorbar=False, cmap="Blues")
plt.tight_layout(); plt.savefig("cm_rf.png"); plt.close()
\end{lstlisting}

    \begin{alertblock}{Teaching note: SITTING vs STANDING confusion}
        Most confused pair: $\sim$13\% inter-confusion. Both are static postures where gravity (9.8 m/s$^2$) dominates. The discriminating signal is subtle gyroscope postural sway. The scatter plot shows \texttt{tBodyAccMag-mean()} separates dynamic vs static activities well (WALKING clusters at $\sim-$0.17, LAYING at $\sim-$0.94); SITTING and STANDING overlap. Feature choice determines what the model can distinguish.
    \end{alertblock}
\end{frame}

%==============================================================================
\section{Exercise 2: Algorithm Comparison}
%==============================================================================

\begin{frame}{Exercise 2: Three Methods --- Same Data}
    \begin{block}{Scenario}
        Before deploying to a device, choose which algorithm to use. Three candidates. Compare on accuracy, inference speed, and hardware deployability.
    \end{block}

    \vspace{0.2cm}
    \textbf{Tasks:}
    \begin{enumerate}
        \item Train kNN (k=5), Random Forest (100 trees), LinearSVC on the same subject split
        \item Measure accuracy, training time, inference time per sample, and serialized model size
        \item Which algorithm fits in 256 kB MCU flash?
        \item Justify choice for: (a) laptop gateway, (b) Raspberry Pi, (c) ARM Cortex-M4 with 256 kB flash
    \end{enumerate}

    \vspace{0.1cm}
    \textbf{Expected conclusion:} SVM wins on accuracy (0.963) and is the only MCU-deployable option (27 kB model). kNN serialized size is $\sim$16 MB (stores all training data). Random Forest is $\sim$5.7 MB (100 trees).
\end{frame}

\begin{frame}[fragile]{Exercise 2: Solution --- Comparison Code}

\begin{lstlisting}[language=Python]
import time, os, pickle, tempfile, pandas as pd
from sklearn.ensemble import RandomForestClassifier
from sklearn.neighbors import KNeighborsClassifier
from sklearn.svm import LinearSVC
from sklearn.preprocessing import StandardScaler
from sklearn.metrics import accuracy_score

X_train = pd.read_csv("har_X_train.csv")
X_test  = pd.read_csv("har_X_test.csv")
y_train = pd.read_csv("har_y_train.csv").squeeze()
y_test  = pd.read_csv("har_y_test.csv").squeeze()

scaler = StandardScaler().fit(X_train)
Xtr_sc = scaler.transform(X_train);  Xte_sc = scaler.transform(X_test)

models = {
    "k-NN (k=5)":    (KNeighborsClassifier(n_neighbors=5, n_jobs=-1), True),
    "Random Forest": (RandomForestClassifier(n_estimators=100, random_state=42, n_jobs=-1), False),
    "SVM (linear)":  (LinearSVC(max_iter=2000, random_state=42), True),
}
for name, (clf, use_scale) in models.items():
    Xtr = Xtr_sc if use_scale else X_train
    Xte = Xte_sc if use_scale else X_test
    t0 = time.time();  clf.fit(Xtr, y_train);  train_t = time.time()-t0
    t0 = time.time();  preds = clf.predict(Xte)
    infer_ms = (time.time()-t0)/len(Xte)*1000
    acc = accuracy_score(y_test, preds)
    with tempfile.NamedTemporaryFile(delete=False) as f:
        pickle.dump(clf, f);  size_kb = os.path.getsize(f.name)/1024;  os.unlink(f.name)
    print(f"{name:<22}  acc={acc:.3f}  train={train_t:.1f}s"
          f"  infer={infer_ms:.3f}ms/sample  size={size_kb:.0f}kB")
\end{lstlisting}
\end{frame}

\begin{frame}{Exercise 2: Expected Results and Teaching Notes}
    \begin{center}
    \begin{tabular}{lccccc}
        \toprule
        \textbf{Method} & \textbf{Accuracy} & \textbf{Train} & \textbf{Infer/sample} & \textbf{Size} & \textbf{MCU?} \\
        \midrule
        k-NN (k=5)    & 0.884 & 0.1 s & $\sim$8 ms    & $\sim$16 MB & No \\
        Random Forest & 0.929 & $\sim$12 s  & $\sim$0.05 ms & $\sim$5.7 MB  & No \\
        SVM (linear)  & \textbf{0.963} & $\sim$4 s & \textbf{$\sim$0.01 ms} & \textbf{27 kB} & \textbf{Yes} \\
        \bottomrule
    \end{tabular}
    \end{center}

    \vspace{0.15cm}
    \fcolorbox{teachborder}{teachbg}{\begin{minipage}{0.93\textwidth}
    \textbf{Teaching notes:}\\
    \textbf{Why SVM wins:} the 561 features are carefully engineered to be linearly separable by the dataset authors. SVM exploits this directly. RF would win on raw time-series where non-linear patterns matter more.\\[4pt]
    \textbf{Why kNN inference is slow:} every prediction computes distances to all 7,352 training samples $\times$ 561 features. At 50 Hz that leaves only 20 ms per sample --- kNN uses most of that budget on a fast laptop and far exceeds it on embedded hardware.\\[4pt]
    \textbf{SVM model size:} one weight matrix (6 classes $\times$ 561 features) = 3,366 float32 values $\approx$ 13 kB raw. Serialized with pickle overhead: $\sim$27 kB measured. Comfortably fits in 256 kB MCU flash.
    \end{minipage}}
\end{frame}

%==============================================================================
\section{Exercise 3: Concept Drift}
%==============================================================================

\begin{frame}{Exercise 3: Deployment Drift}
    \begin{block}{Scenario}
        The hospital deploys the SVM to 500 patients including elderly adults. After 6 months, accuracy drops. Elderly gait has lower amplitude, phone placement varies, accelerometer offset has drifted.
    \end{block}

    \vspace{0.2cm}
    \textbf{Drift simulation:} scale all 345 acceleration feature columns (those with ``Acc'' in the name) by 0.75 and add Gaussian noise. Grounded in Menz et al.\ (2003): elderly adults produce 10--20\% lower gait acceleration amplitude. Scale 0.75 gives acc $\approx$0.74 --- a realistically severe but plausible deployment shift.

    \vspace{0.1cm}
    \textbf{Tasks:}
    \begin{enumerate}
        \item Measure SVM accuracy on the drifted test set (scale 0.75)
        \item Plot accuracy vs drift magnitude (scale 1.0 down to 0.70)
        \item Print per-class recall on the drifted set. Which activities degrade most?
        \item Retrain on mixed data (original + 40\% drifted re-labeled). Measure recovery.
    \end{enumerate}
\end{frame}

\begin{frame}[fragile]{Exercise 3: Solution --- Drift Simulation}

\begin{lstlisting}[language=Python]
import numpy as np, pandas as pd, matplotlib.pyplot as plt
from sklearn.svm import LinearSVC
from sklearn.preprocessing import StandardScaler
from sklearn.metrics import accuracy_score

X_train = pd.read_csv("har_X_train.csv")
X_test  = pd.read_csv("har_X_test.csv")
y_train = pd.read_csv("har_y_train.csv").squeeze()
y_test  = pd.read_csv("har_y_test.csv").squeeze()

scaler = StandardScaler().fit(X_train)
svm    = LinearSVC(max_iter=2000, random_state=42)
svm.fit(scaler.transform(X_train), y_train)
baseline = accuracy_score(y_test, svm.predict(scaler.transform(X_test)))
print(f"Baseline (young adults): {baseline:.3f}")

acc_cols = [c for c in X_train.columns if "Acc" in c]  # 345 cols
rng = np.random.default_rng(7)
results = []
for scale in np.arange(1.0, 0.69, -0.05):
    Xd = X_test.copy()
    Xd[acc_cols] = Xd[acc_cols] * scale
    Xd = Xd + rng.normal(0, 0.03*(1.2-scale), Xd.shape)
    acc = accuracy_score(y_test, svm.predict(scaler.transform(Xd)))
    results.append({"scale": round(scale,2), "accuracy": acc})

df_r = pd.DataFrame(results)
plt.figure(figsize=(8,4))
plt.plot(df_r["scale"], df_r["accuracy"], "o-r", label="Drifted accuracy")
plt.axhline(baseline, linestyle="--", color="blue",
            label="Baseline %.3f" % baseline)
plt.axhline(0.85, linestyle=":", color="orange", label="Alert threshold 0.85")
plt.gca().invert_xaxis()
plt.xlabel("Acc feature scale (1.0=young adults, 0.75=elderly gait)")
plt.ylabel("Accuracy");  plt.title("SVM Accuracy Under Gait Drift")
plt.legend();  plt.tight_layout();  plt.savefig("drift_plot.png");  plt.close()
\end{lstlisting}
\end{frame}

\begin{frame}[fragile]{Exercise 3: Solution --- Retraining}

\begin{lstlisting}[language=Python]
# Mixed retraining: original training + 40% of drifted test re-labeled
rng      = np.random.default_rng(7)
acc_cols = [c for c in X_train.columns if "Acc" in c]  # 345 cols

X_new = X_test.copy().sample(frac=0.4, random_state=1)
y_new = y_test.loc[X_new.index]
X_new[acc_cols] = X_new[acc_cols] * 0.75
X_new = X_new + rng.normal(0, 0.03*(1.2-0.75), X_new.shape)

X_mixed = pd.concat([X_train, X_new], ignore_index=True)
y_mixed = pd.concat([y_train, y_new], ignore_index=True)

scaler_v2 = StandardScaler().fit(X_mixed)
svm_v2    = LinearSVC(max_iter=2000, random_state=42)
svm_v2.fit(scaler_v2.transform(X_mixed), y_mixed)

X_drifted = X_test.copy()
X_drifted[acc_cols] = X_drifted[acc_cols] * 0.75
X_drifted = X_drifted + rng.normal(0, 0.03*(1.2-0.75), X_drifted.shape)

acc_v1 = accuracy_score(y_test, svm.predict(scaler.transform(X_drifted)))
acc_v2 = accuracy_score(y_test, svm_v2.predict(scaler_v2.transform(X_drifted)))
acc_v2_clean = accuracy_score(y_test, svm_v2.predict(scaler_v2.transform(X_test)))
print(f"v1 on drifted (elderly): {acc_v1:.3f}")
print(f"v2 on drifted (elderly): {acc_v2:.3f}")
print(f"v2 on clean  (young):   {acc_v2_clean:.3f}")
\end{lstlisting}

    \begin{exampleblock}{Expected output (measured)}
        v1 on drifted: \textbf{0.742}. v2 on drifted: \textbf{0.976}. v2 on clean: \textbf{0.977}. Full recovery because the 40\% new training samples come from the same test subjects (best-case). Real deployments with genuinely new patients will show partial but meaningful recovery. The direction is always correct: mixed retraining improves on the drifted distribution.
    \end{exampleblock}
\end{frame}

\begin{frame}[fragile]{Exercise 3: Per-Activity Drift --- Teaching Notes}
    \fcolorbox{teachborder}{teachbg}{\begin{minipage}{0.93\textwidth}
    \textbf{Key insight: drift is class-specific, not uniform.}\\[4pt]
    \textbf{Activities that degrade most (measured at scale 0.75):} WALKING (recall 1.000$\to$0.379, $-$0.621) and WALKING\_UPSTAIRS (1.000$\to$0.306, $-$0.694). SITTING drops moderately (0.821). These are acceleration-dependent activities. Their discriminating features are directly scaled by the 0.75 factor.\\[4pt]
    \textbf{Activities largely unaffected:} SITTING, STANDING, LAYING. Both are static --- acceleration dominated by gravity ($\sim$9.8 m/s$^2$), which does not change with gait amplitude.\\[4pt]
    \textbf{Implication for monitoring:} average accuracy can look ``ok'' while walking classification completely fails. Class-level monitoring is necessary. Have students print \texttt{classification\_report} on the drifted test set to observe this directly.
    \end{minipage}}

    \vspace{0.2cm}
\begin{lstlisting}[language=Python]
from sklearn.metrics import classification_report
print("--- SVM v1 on drifted test ---")
print(classification_report(y_test, svm.predict(scaler.transform(X_drifted))))
\end{lstlisting}
\end{frame}

%==============================================================================
\section{Take-Home: Solutions and Rubric}
%==============================================================================

\begin{frame}[fragile]{Take-Home Task 1 Solution: Confusion Matrix + Feature Importance}

\begin{lstlisting}[language=Python]
import numpy as np, matplotlib.pyplot as plt
from sklearn.metrics import ConfusionMatrixDisplay

ConfusionMatrixDisplay.from_estimator(
    svm, scaler.transform(X_test), y_test,
    normalize="true", xticks_rotation=45, cmap="Blues", colorbar=False)
plt.tight_layout(); plt.savefig("cm_svm_norm.png"); plt.close()

coef_imp = np.abs(svm.coef_).mean(axis=0)
top15    = np.argsort(coef_imp)[::-1][:15]
print("Top 15 features by SVM weight:")
for i in top15:
    print(f"  {X_test.columns[i]:<42}  {coef_imp[i]:.4f}")
\end{lstlisting}

    \begin{exampleblock}{Expected findings}
        \textbf{Most confused pair:} SITTING vs STANDING ($\sim$12\% inter-confusion). Both static; discriminating signal is gyroscope postural sway.\\[4pt]
        \textbf{Top features (measured):} gyroscope jerk entropy (\texttt{tBodyGyroJerk-entropy()-X}) has the highest SVM weight, followed by acceleration jerk entropy. Entropy features capture irregularity of motion, which is very discriminative. Frequency-domain features appear in the top 15 but time-domain entropy dominates. Key insight: the most important features are not the means --- they are the distributional shape descriptors.
    \end{exampleblock}
\end{frame}

\begin{frame}[fragile]{Take-Home Task 2 Solution: KS-Test Drift Detection}

\begin{lstlisting}[language=Python]
import numpy as np, pandas as pd, matplotlib.pyplot as plt
from scipy.stats import ks_2samp

rng      = np.random.default_rng(7)
acc_cols = [c for c in X_test.columns if "Acc" in c]  # 345 cols
X_drifted = X_test.copy()
X_drifted[acc_cols] = X_drifted[acc_cols] * 0.75
X_drifted = X_drifted + rng.normal(0, 0.03*(1.2-0.75), X_drifted.shape)

rows = []
for col in X_test.columns:
    stat, p = ks_2samp(X_test[col].values, X_drifted[col].values)
    rows.append({"feature": col, "KS": stat, "p": p,
                 "drifted": p < 0.05, "is_acc": "Acc" in col})
df_ks = pd.DataFrame(rows).sort_values("KS", ascending=False)
print(f"Features drifted (p<0.05): {df_ks['drifted'].sum()} / {len(df_ks)}")
print(f"  Acc: {df_ks[df_ks.is_acc & df_ks.drifted].shape[0]}"
      f"/{df_ks.is_acc.sum()}  "
      f"  Gyro: {df_ks[~df_ks.is_acc & df_ks.drifted].shape[0]}"
      f"/{(~df_ks.is_acc).sum()}")

top20  = df_ks.head(20)
colors = ["red" if d else "steelblue" for d in top20["drifted"]]
plt.figure(figsize=(12,4))
plt.bar(range(20), top20["KS"], color=colors)
plt.xticks(range(20), [f[:20] for f in top20["feature"]],
           rotation=45, ha="right", fontsize=7)
plt.title("Top 20 features by KS statistic (red = p<0.05)")
plt.ylabel("KS statistic")
plt.tight_layout(); plt.savefig("ks_drift.png"); plt.close()
\end{lstlisting}
\end{frame}

\begin{frame}{Take-Home Task 2: Expected Results and Teaching Notes}
    \textbf{Expected output (measured on real data):}
    \begin{itemize}
        \item 344 out of 345 acceleration features flagged (p $<$ 0.05, 99.7\%)
        \item 148 out of 213 gyroscope features also flagged (69\%) --- because gyroscope and accelerometer features are correlated (body motion affects both sensors)
        \item Total: 492 / 561 features drifted --- KS test fires a clear global alarm
        \item The \textbf{acceleration flagging rate is higher} (99.7\% vs 69\%), correctly identifying the primary drift source
    \end{itemize}

    \vspace{0.2cm}
    \fcolorbox{teachborder}{teachbg}{\begin{minipage}{0.93\textwidth}
    \textbf{Key discussion point:}\\
    The KS test detected drift \textbf{without any labels}. In production, after deploying to elderly patients, ground truth is not available immediately --- confirming the activity requires a nurse to watch video, days later. KS tests run on incoming feature distributions \textit{the moment data arrives}. This fires an early warning \textbf{days before} accuracy-based monitoring can confirm the problem.\\[4pt]
    \textbf{This is the production value of unsupervised drift detection.}\\[4pt]
    The correct monitoring strategy depends on label availability: if you have continuous ground truth (e.g., supervised follow-up), accuracy monitoring is better. If labels are delayed or expensive, KS-based monitoring is the only real-time option.
    \end{minipage}}
\end{frame}

\begin{frame}{Take-Home Task 3: Grading Rubric}
    \begin{center}
    \footnotesize
    \begin{tabular}{p{3.0cm}p{3.8cm}p{3.5cm}c}
        \toprule
        \textbf{Criterion} & \textbf{Full marks} & \textbf{Partial marks} & \textbf{Wt.} \\
        \midrule
        Sensor + features & 4--8 physically motivated features; extraction window and rate stated & Features listed without physical justification & 20\% \\
        Training data & Classes defined; labeling method realistic for IoT; bottleneck identified & Classes defined but labeling strategy vague & 20\% \\
        Algorithm choice & Correct for hardware target; justified with size and inference time; tradeoff stated & Algorithm named without hardware justification & 25\% \\
        Drift causes & 2+ physically specific causes for this sensor type & ``Data will change'' (too generic) & 20\% \\
        Monitoring & Accuracy vs KS choice justified by label availability & Monitoring mentioned without justification & 15\% \\
        \bottomrule
    \end{tabular}
    \normalsize
    \end{center}

    \vspace{0.1cm}
    \begin{alertblock}{Common weak answers to push back on}
        \textbf{Algorithm:} ``I would use Random Forest because it is accurate.'' Expected: ``I need $<$50 kB for Cortex-M4; RF at 5.7 MB is too large; SVM serialized size $\approx$27 kB fits.''\\
        \textbf{Drift:} ``sensor data changes over time.'' Expected: ``the MEMS accelerometer baseline offset drifts 0.05--0.1 g/year due to mechanical stress on the proof mass suspension.''
    \end{alertblock}
\end{frame}

%==============================================================================
\section{Reference}
%==============================================================================

\begin{frame}{Requirements and Setup Reference}
    \begin{columns}
        \column{0.48\textwidth}
        \textbf{Python dependencies:}
        \begin{tabular}{ll}
            \toprule
            \textbf{Package} & \textbf{Purpose} \\
            \midrule
            \texttt{ucimlrepo>=0.0.7} & dataset download \\
            \texttt{scikit-learn>=1.3} & all ML methods \\
            \texttt{pandas>=2.0}      & data handling \\
            \texttt{numpy>=1.24}      & arrays \\
            \texttt{matplotlib>=3.7}  & plots \\
            \texttt{seaborn>=0.12}    & heatmaps \\
            \texttt{scipy>=1.11}      & KS test \\
            \bottomrule
        \end{tabular}

        \vspace{0.2cm}
        \textbf{Install:}\\
        {\small \texttt{pip install ucimlrepo scikit-learn}}\\
        {\small \texttt{pandas numpy matplotlib seaborn scipy}}

        \column{0.48\textwidth}
        \textbf{Dataset reference:}\\
        Reyes-Ortiz, Anguita, Ghio, Oneto, Parra (2013). \textit{HAR Using Smartphones}. UCI ML Repository. \texttt{doi:10.24432/C54S4K}

        \vspace{0.2cm}
        \textbf{Drift reference:}\\
        Menz, Lord, Fitzpatrick (2003). Age-related differences in walking stability. \textit{Age and Ageing}, 32(2), 137--142.

        \vspace{0.2cm}
        \textbf{Key scikit-learn classes:}\\
        \texttt{sklearn.svm.LinearSVC}\\
        \texttt{sklearn.ensemble.RandomForestClassifier}\\
        \texttt{sklearn.neighbors.KNeighborsClassifier}\\
        \texttt{sklearn.metrics.ConfusionMatrixDisplay}\\
        \texttt{scipy.stats.ks\_2samp}
    \end{columns}
\end{frame}

\begin{frame}{Questions?}
    \begin{center}
        \Huge Questions?

        \vspace{0.8cm}
        \normalsize
        \textbf{Eduardo Baena}\\
        \texttt{e.baena@northeastern.edu}

        \vspace{0.5cm}
        \textit{EECE5155: Wireless Sensor Networks and the Internet of Things}\\
        Spring 2026

        \vspace{0.4cm}
        \textbf{Pipeline:} Real sensor data $\rightarrow$ Feature engineering $\rightarrow$ Train $\rightarrow$
        Compare methods $\rightarrow$ Deploy $\rightarrow$ Drift $\rightarrow$ Retrain
    \end{center}
\end{frame}

\end{document}
